\section{Work Objective}

This thesis aims to improve the accuracy of Estimated Time of Arrival (ETA) prediction by developing a novel dynamic Graph Neural Network (GNN) architecture that models traffic at the vehicle level with explicit route information. Unlike existing approaches that either use static sensor networks (e.g., DCRNN~\cite{dcrnn2018}, ST-GCN~\cite{yu2018stgcn}) or infer routes implicitly from origin-destination pairs (e.g., DeepTTE~\cite{deepTTE2018}, STAD~\cite{abbar2020stad}), this work proposes a hybrid graph structure where static infrastructure (junctions as nodes, roads as edges) is combined with dynamic vehicle nodes and time-varying edges. The dynamic edges create junction-vehicle-vehicle-junction relations by connecting vehicles to junctions (traversal edges) and vehicles to other vehicles on the same road segment (interaction edges), enabling rich modeling of vehicle interactions and traffic flow dynamics.

The primary research objective is to investigate whether incorporating predefined vehicle routes as explicit route information significantly enhances ETA prediction accuracy compared to models that infer routes implicitly (e.g., DeepTTE~\cite{deepTTE2018}, STAD~\cite{abbar2020stad}). Current state-of-the-art methods such as DuETA~\cite{dueta2023} use route information but focus on aggregated road segments rather than individual vehicle trajectories. This work will develop and evaluate a dynamic GNN architecture that explicitly encodes route intent at the vehicle level, comparing its performance against baseline models that infer routes implicitly (e.g., DeepTTE~\cite{deepTTE2018}, STAD~\cite{abbar2020stad}) or use static sensor networks (e.g., DCRNN~\cite{dcrnn2018}, ST-GCN~\cite{yu2018stgcn}). The research will demonstrate the contribution of route-aware modeling combined with vehicle-level dynamic graph representation to achieving more accurate ETA predictions in dynamic traffic scenarios.

